% !PRINTED 376
% No folio or running head is printed on this chapter-opening leaf; printed page 376 is established by adjacent numbered leaves.
\weberpage{376}
\begin{center}
\gesperrt{Fifteenth Chapter.}

Genera of the Quadratic Forms.
\end{center}
\begin{center}\rule{0.2\linewidth}{0.4pt}\end{center}

\begin{center}
\textbf{\S\,103. \gesperrt{Representation of Numbers by Quadratic Forms.}}
\end{center}

A natural number $m$ is called \gesperrt{properly representable} by the primitive form $(a,b,c)$ with the discriminant $D$ if there are two relatively prime numbers $x$, $y$ which satisfy the equation
\weq{1}{ax^2 + bxy + cy^2 = m}
and one can then find a form $(m,n,l)$ equivalent to $(a,b,c)$ whose discriminant
\weq{2}{D = n^2 - 4ml}
% [sic] The source uses the singular “this class” despite naming two classes; the wording is retained.
is. If $m$ is representable by the form $(a,b,c)$, then it is also representable by every equivalent form, and hence by every form of the class $A$ to which $(a,b,c)$ belongs. We then call $m$ representable by the \gesperrt{class $A$}. If $A$, $A'$ are two classes and $m$, $m'$ two mutually coprime numbers which are properly representable by this class, then the product $mm'$ is properly representable by the composed class $AA'$. This follows from the compositions of the forms (12), (14), \S\,101.

From (2) there resulted:
\weq{3}{n^2 \equiv D \ (\mathrm{mod}\ 4m).}

If $(m,n,l)$ is imprimitive, then $D$ cannot be a fundamental discriminant. As hitherto, we denote the fundamental discriminant associated with $D$ by $\Delta$ and put
\[
D = Q^2 \Delta,
\]
and the greatest common divisor of $m$, $n$, $l$ divides $Q^2$. In what follows we therefore assume that $m$ \gesperrt{is relatively prime} to $Q$, and are then certain that the \gesperrt{form $(m,n,l)$ determined by (2) is primitive}.
